\begin{center}
    {\LARGE \textbf{2010物理化学}} \\[2ex]
    {\large 时量180分钟 \quad 满分90分}
\end{center}

\vspace{0.5cm}
\noindent\textbf{注意事项:}
\begin{enumerate}[label=\arabic*., parsep=0pt, itemsep=0pt]
    \item 答题前,考生先将自己的姓名、学号、考场号、座位号填写在答题卡上,将条形码准确粘贴在考生信息条形码粘贴区。
    \item 考试结束后,将本试卷与答题纸一并收回。
    \item 回答各个问题时,将答案写在答题纸上,写在本试卷上无效。
\end{enumerate}

\vspace{0.5cm}
\noindent\textbf{1. (15 pts.)}

在两个中间由旋塞连通的烧瓶中,开始时分别盛有$0.2\ \mathrm{mol},0.2p^\ominus$的$\ce{O_{2}}$和$0.8\ \mathrm{mol},0.8p^\ominus$

的$\ce{N_{2}}$(均视为理想气体),置于$25^\circ\mathrm{C}$的恒温水浴里,然后打开旋塞。
\begin{enumerate}[label=(\arabic*), itemsep=1ex]
    \item 试计算最终$p_2$；
    \item 计算该过程的$Q,W,\Delta_\mathrm{mix}H,\Delta_\mathrm{mix}S$和$\Delta_\mathrm{mix}G$。
\end{enumerate}

\vspace{0.5cm}
\noindent\textbf{2. (15 pts.)}

$300\ \mathrm{K}$时,试计算(1) 从大量的等物质的量的$\ce{C_{2}H_{4}Br_{2}}$和$\ce{C_{3}H_{6}Br_{2}}$的理想混合物中分离出$1\ \mathrm{mol}$纯$\ce{C_{2}H_{4}Br_{2}}$过程的$\Delta G_1$,(2) 若混合物中各含$2\ \mathrm{mol}$ $\ce{C_{2}H_{4}Br_{2}}$和$\ce{C_{3}H_{6}Br_{2}}$,从中分离出$1\ \mathrm{mol}$纯$\ce{C_{2}H_{4}Br_{2}}$时,$\Delta G_2$又为多少。

\vspace{0.5cm}
\noindent\textbf{3. (15 pts.)}

市售民用高压锅内的压力可达$233\ \mathrm{kPa}$,问此时水的沸点为多少。已知水的汽化热为$40.67\ \mathrm{kJ\cdot mol^{-1}}$。

\vspace{0.5cm}
\noindent\textbf{4. (15 pts.)}

\begin{enumerate}[label=(\arabic*), itemsep=1ex]
    \item $\ce{2NaHCO_{3}(s) = Na_{2}CO_{3}(s) + H_{2}O(g) + CO_{2}(g)}$
    \item $\ce{CuSO_{4} * 5H_{2}O(s) = CuSO_{4} * 3H_{2}O(s) + 2H_{2}O(g)}$
\end{enumerate}
反应(1)和(2)在$50^\circ\mathrm{C}$时的分解压力分别为$p(1)=4.00\ \mathrm{kPa}$,$p(2)=6.05\ \mathrm{kPa}$。若分解反应(1)和(2)在同一容器内进行,计算$50^\circ\mathrm{C}$时体系的平衡压力。

\vspace{0.5cm}
\noindent\textbf{5. (15 pts.)}

某一特定的一级反应在$27^\circ\mathrm{C}$反应时,经过$5000\ \mathrm{s}$后,反应物的浓度减少到初始值的一半,在$37^\circ\mathrm{C}$时,经过$1000\ \mathrm{s}$,浓度就减半,计算:
\begin{enumerate}[label=(\arabic*), itemsep=1ex]
    \item $27^\circ\mathrm{C}$时的反应速率常数；
    \item 在$37^\circ\mathrm{C}$反应时,当反应物浓度降低到其初始值的四分之一时所需的时间；
    \item 该反应的活化能。
\end{enumerate}

\vspace{0.5cm}
\noindent\textbf{6. (15 pts.)}

有电池$\ce{Ag|AgCl(s)|HCl(0.1\ mol\cdot kg^{-1})|Cl_{2}(p^\ominus)|Pt}$,

已知$\ce{AgCl}$在$25^\circ\mathrm{C}$时的标准生成焓为$-127.03\ \mathrm{kJ\cdot mol^{-1}}$,$\ce{Ag}$、$\ce{AgCl}$和$\ce{Cl_{2}(g)}$在$25^\circ\mathrm{C}$时的标准熵依次为$41.95\ \mathrm{J\cdot K^{-1}\cdot mol^{-1}}$,$96.10\ \mathrm{J\cdot K^{-1}\cdot mol^{-1}}$和$243.86\ \mathrm{J\cdot K^{-1}\cdot mol^{-1}}$。

试计算$25^\circ\mathrm{C}$时,
\begin{enumerate}[label=(\arabic*), itemsep=1ex]
    \item 电池电动势；
    \item 电池可逆操作时的热效应；
    \item 电池的温度系数。
\end{enumerate}